\section{Relational E-Matching}
\label{sec:algorithm}

% {\centering
% \begin{minipage}{.8\linewidth}
% \begin{algorithm}%[H]
%   \LinesNumbered{} \SetAlgoLined{} \KwIn{a function symbol $f$ and variable
%     ordering $\pi$. } \KwOut{ a trie $R_{f} : \{ \pi((c, \bar x)) \mid f(\bar
%     x) \in G[c] \}$. } \For{ $f(\bar x) \in G$ }{ $R_{f}$.insert$(\pi(f.\eclass,
%     \bar x))$ }
%   \caption{ Converting $f$-nodes in $G$ into a relation (as a trie).
%     \yz{should we keep this?}
%   }\label{alg:g2trie}
% \end{algorithm}
% \end{minipage}
% \par}

\Ematching via
 backtracking search is inefficient because it handles equality constraints
 suboptimally.
In fact, backtracking search follows edges in the \egraph and only
 visits concrete terms that satisfy structural constraints.
However, equality constraints
 are checked \emph{a posteriori} only after the search visits a (partial) term.\footnote{
   Backtracking \ematching can perform the check as soon as it has traversed
   enough of the pattern to encounter a variable more than once.
}
Whenever there are many
 terms that satisfy the structural constraints but not the equality constraints,
 as is in our example pattern $f(\alpha, g(\alpha))$,
 backtracking will waste time visiting terms that do not yield a match.

By reducing \ematching to evaluating conjunctive queries,
 we can use join algorithms that take advantage of both structural
 \emph{and} equality constraints.
\autoref{fig:emvhj} conveys this intuition
 using the pattern $f(\alpha, g(\alpha))$ and the
 example \egraph and database from \autoref{fig:egraph-and-table}.
The backtracking approach considers every possible assignment
 to the variables, even those where the two occurrences of $\alpha$
 do not agree.

We can instead formulate a conjunctive query
 that is equivalent to the following pattern:
 $$Q(root, \alpha) \gets R_f(root, \alpha, x), R_g(x, \alpha).$$
Later subsections will detail how this conversion is done,
 but note how the auxiliary variable $x$ captures the structural
 constraint from the pattern.
Evaluating $Q$ with a simple hash join strategy exemplifies the
 benefits of the relational approach:
 it considers structural \emph{and} equality constraints
 (in this case by doing a hash join keyed on $(x, \alpha)$);
 indeed, the relational perspective sees
 no difference between the two kinds of constraints.

\begin{figure}
  \begin{subfigure}{0.45\textwidth}
    \centering
    \begin{tabular}{ c c c }
      $f(1, i_{g})$ & $g(1)$ & \checkmark\\
      \multicolumn{3}{c}{\vdots} \\
                    & $g(N)$ & $\times$\\
      $f(2, i_{g})$ & $g(1)$ & $\times$\\
                    & $g(2)$ & \checkmark\\
      \multicolumn{3}{c}{\vdots} \\
                    & $g(N)$ & $\times$\\
      $f(N, i_{g})$ & $g(1)$ & $\times$\\
      \multicolumn{3}{c}{\vdots} \\
                    & $g(N)$ & \checkmark\\
    \end{tabular}
    \caption{ Backtracking takes time $O(N^{2})$ }\label{tab:hj}
  \end{subfigure}
  \begin{subfigure}{0.45\textwidth}
    \centering $\rotatebox[origin=c]{90}{build hash}%
    \left\downarrow
      \begin{tabular}{ c }
        $R_g(i_g, 1)$ \\
        $R_g(i_g, 2)$ \\
        $R_g(i_g, 3)$ \\
        $\vdots$ \\
        $R_g(i_g, N)$ \\
      \end{tabular}
    \right.$ $\rotatebox[origin=c]{90}{probe}%
    \left\downarrow
      \begin{tabular}{ c c }
        $R_f(i_f, 1, i_g)$ & \checkmark \\
        $R_f(i_f, 2, i_g)$ & \checkmark\\
        $R_f(i_f, 3, i_g)$ & \checkmark\\
        $\vdots$ & \\
        $R_f(i_f, N, i_g)$ & \checkmark\\
      \end{tabular}
    \right.$
    \caption{ Hash join takes time $O(N)$. }\label{tab:hj}
  \end{subfigure}
  \caption{
    \Ematching $f(\alpha, g(\alpha))$
    with backtracking search and a simple hash join
    on the \egraph/database in \autoref{fig:egraph-and-table}.
  }\label{fig:emvhj}
\end{figure}


This observation leads us to a very simple
 algorithm for relational \ematching,
 shown in \autoref{alg:main}.
Relational \ematching~takes an e-graph $E$ and a set
 of patterns $\textit{ps}$.
It first transforms the \egraph to a relational database $I$.
Then, it reduces every pattern $p$ to a conjunctive query $q$.
Finally, it evaluates the conjunctive queries over $I$.
These intermediate steps will be detailed in the following subsections.

\begin{algorithm}
  \LinesNumbered{} \SetAlgoLined{}
  \KwIn{An \egraph~$E$ and a list of \ematching~patterns $\textit{ps}$}
  \KwOut{The result of running $\textit{ps}$ on $E$}
  $I  \gets \textsc{EGraphToDatabase}(E)$\;
  $qs \gets \{ \textsc{PatternToCQ}(p) \mid p \in ps \}$\;
  \Return $\{\textsc{EvalCQ}(q, I)\ \mid q \in qs \}$
  \caption{\textsc{RelationalEMatching}}
  \label{alg:main}
\end{algorithm}

\subsection{From the E-Graph to a Relational Database}

% \begin{figure}
%   \centering
%   \begin{align*}
%     \textsc{EgraphToDatabase}(G) &= \{\textsc{EnodeToAtom}(n)\ \mid\text{ for \enode\ } n\in G\}\\
%     \textsc{EnodeToAtom}(f(c_1,\ldots,c_k)) &= R_f(\textit{lookup}(f(c_1,\ldots,c_k)), c_1,\ldots, c_k)
%   \end{align*}
%   \caption{Transforming an \egraph~to relational database}
%   \label{fig:egraph_to_db}
% \end{figure}

\begin{figure}
  \begin{subfigure}{0.45\textwidth}
    \centering \includegraphics[width=0.65\linewidth]{fgn.pdf}
    \caption{An example \egraph, reproduced from \autoref{fig:egraph-eg}.}\label{fig:egraph-eg2}
  \end{subfigure}
  \hfill
  \begin{subfigure}{0.25\textwidth}
    \centering
    \begin{tabular}{ c c c }
      \toprule
      id & $\text{arg}_{1}$ & $\text{arg}_{2}$ \\
      \midrule
      $i_{f}$ & 1 & $i_{g}$ \\
      $i_{f}$ & 2 & $i_{g}$ \\
      \vdots & \vdots & \vdots \\
      $i_{f}$ & $N$ & $i_{g}$ \\
      \bottomrule
    \end{tabular}
    \caption{Relation of $f$.}\label{tab:rel-f}
  \end{subfigure}
  \hfill
  \begin{subfigure}{0.25\textwidth}
    \centering
    \begin{tabular}{ c c }
      \toprule
      id & $\text{arg}_{1}$ \\
      \midrule
      $i_{g}$ & 1 \\
      $i_{g}$ & 2 \\
      \vdots & \vdots \\
      $i_{g}$ & $N$ \\
      \bottomrule
    \end{tabular}
    \caption{Relation of $g$.}\label{tab:rel-g}
  \end{subfigure}
  \caption{ An \egraph~and its relational representation. Each \eclass~(dotted
    box) is labeled with its id.
  }\label{fig:egraph-and-table}
\end{figure}

The first step of relational e-matching is to transform the \egraph $E$
 into a relational database $I$.
The domain of the database is \eclass ids,
 and its schema is determined
 by the function symbols in $\Sigma$.
Every \enode with symbol $f$ in the \egraph
 corresponds to a tuple in the relation $R_f$ in the database.
If $f$ has arity $k$,
 then $R_f$ will have arity $k + 1$;
 its first attribute is the \eclass id that contains
 the corresponding \enode,
 and the remaining attributes are the $k$ children of the $f$ \enode.
\autoref{fig:egraph-and-table} shows an example \egraph
 and part of its corresponding database.
In particular, only the relations
 of function symbols $f$ and $g$ are presented in this figure.
There are $N$ other relations;
 each relation $R_j$ represents a constant
 $j$ and has exactly one tuple (i.e., singleton $(j)$).

We construct the database $I$
 by simply looping over every \enode in every \eclass in $E$
 and making
 a tuple in the corresponding relation:
\[
  I =
  \left\{
    R_f \gets (\find(i), \find(j_1), \ldots, \find(j_k))
    \mid
    M[i] = f(j_1, \ldots, j_k)
  \right\}
\]
Note that the tuples in the database contain only canonical \eclass
 ids returned from the \find{} function.

% An \egraph~$E$ representing a congruence relation over
% $T(\Sigma,\emptyset)$ will be mapped to a relational database over $C$, the set
% of \eclasses, with schema $S_C$, and every \enode~in $E$ will correspond to an
% atom in this relational database. Specifically, for every $k$-ary symbol $f$ in
% $\Sigma$, We add an $k+1$-ary symbol $R_f$ to $S_C$. We refer to fields of $R_f$
% as $\eclass, \text{arg}_1,\ldots,\text{arg}_k$ respectively, and they denote the
% \eclass~ an \enode~belongs to and the children \eclasses~of this
% \enode~respectively. For every \enode~in the \egraph, relational
% \ematching~insert an atom in the relational database. The algorithm is presented
% in \autoref{fig:egraph_to_db}.

% Take the \egraph~in Figure~\ref{fig:egraph-eg} for example. The transformed
% relation $R_f$ and $R_g$ for function symbol $f$ and $g$ are presented in
% Figure~\ref{tab:rel-f} and Figure~\ref{tab:rel-g}. The $f$-application
% \enodes~in \eclass~$c_f$ translate to atom $R_f(c_f, i,c_g)$ for $i\in[N]$,
% because the \enodes~are contained in \eclass~$c_f$ and have children $i$ and
% $c_g$ for some $i\in[N]$. Similarly, the $g$-application \enode~in \eclass~$c_g$
% translates to the atom $R_g(c_g, i)$ for $i\in[N]$. \autoref{fig:egraph_to_db}
% only shows the relation representing $f$ and $g$. Besides the relation $R_f$ and
% $R_g$ presented, there are $N$ other relations, each representing a constant
% symbol and contains one entry that denotes the \eclass~where the corresponding
% \enode~is contained.

Our presentation in this paper specifically targets
 \ematching use cases like equality saturation,
 where the building of the database $I$ can be amortized.
In this setting, \ematching is done in large batches,
 and expensive work like congruence closure can be
 amortized between these batches using a technique called ``rebuilding''~\cite{egg}.
The time complexity of building this database is always linear,
 which is subsumed by the time complexity
 of most non-trivial \ematching patterns.
In \autoref{sec:incremental},
 we discuss how this technique could be
 generalized to the non-amortized setting of
 frequently updated \egraph as future work.


\subsection{From Patterns to Conjunctive Queries}
\label{sec:pattern-to-query}

\def\aux{\textsc{Aux}} \def\compile{\textsc{Compile}}
\begin{figure}
  \begin{align*}
    \compile(p) &= Q(\textit{root}, v_1,\ldots,v_k) \leftarrow \textit{atoms}\\
                &\text{\quad where $v_1\dots v_k$ are variables in $p$}\\
                &\text{\quad and $\aux(p) = \textit{root} \sim \textit{atoms}$}
    \\
    \aux(f(p_1,\ldots, p_k)) &= v \sim R_f(v, v_1, \ldots, v_k), A_1, \ldots, A_k \\
                & \text{\quad { where} $v$ is fresh and }\aux(p_i) = v_i \sim A_i
    \\
    \aux(x) &= x \sim \emptyset \text{\qquad { where} $x$ is a pattern variable}\\
  \end{align*}
  \caption{Compiling a pattern to a conjunctive query.}\label{alg:compile}
\end{figure}



Once we have a database that corresponds to the \egraph,
 we must convert each pattern we wish to \ematch to a conjunctive query.
We use the algorithm in
\autoref{alg:compile} to ``unnest'' a pattern to a conjunctive query
 by connecting nested patterns with \textit{auxiliary variables}.

The \aux{}
function returns a variable and a conjunctive query atom list.
Particularly, for
non-variable pattern $f(p_1,\ldots p_k)$,
\aux{} produces a fresh variable $v$
and a concatenation of $R_f(v, v_1,\ldots, v_k)$ and atoms from $A_i$, where
$v_i\sim A_i$ is the result of calling $\aux(p_i)$.
For variable pattern $x$,
 \aux{} simply returns $x$ and an empty list.
Note that the auxiliary variables introduced by
 $\aux(f(\ldots))$ are not included in the head of the query,
 and thus are not part of the output.

Given a pattern $p$, the \compile{}
function returns a conjunctive query with body atoms from $\aux(p)$ and the head
atom consisting of the root variable and variables in $p$. The compiled
conjunctive query and the original \ematching~query are equivalent because there
is a one-to-one correspondence between the output of them.
Specifically, each
 \ematching output
 $(i_{\textit{root}}, \sigma)$
 corresponds to a query output of
 $\{ \textit{root} \mapsto i_{\textit{root}} \} \cup \sigma$.
The only difference is that
 returning the root \eclass id is a special consideration for \ematching,
 but it is just another variable in the conjunctive query.

The \compile{} function (specifically the \aux{} subroutine) relies on
 the fact that the database contains only canonical \eclass ids.
Without this fact, nested patterns would require an additional join
 on the equivalence relation $\equivid$.
But since $i \equivid j \iff i = j$
 if $i$ and $j$ are canonical \eclass ids,
 we can omit introducing the additional join,
 instead joining nested patterns directly on the auxiliary variable.

Using this algorithm, the example pattern $f(\alpha, g(\alpha))$ is compiled to
the following conjunctive query:
\begin{equation}
  \label{eqn:nonlinear-cq}
  Q(\mathit{root},\alpha)\gets
  R_f(\mathit{root},\alpha,x),R_g(x,\alpha).
\end{equation}
Compared to the original
\ematching~pattern, this flattened representation enables relational \ematching to
utilize both the structural and the equality constraints. For example, a reasonable
query plan that database optimizers will synthesize is a hash join
on both join variables (i.e., $x$ and $\alpha$), which takes
$O(N)$ time. In contrast, backtracking-based \ematching{} takes $O(N^2)$ time.

Figure~\ref{fig:emvhj} shows the traces for running a direct backtracking search on the
\egraph~and running hash join on the relational representation.
Every term enumerated by hash join will simultaneously satisfy all the constraints.
Conceptually, backtracking-based \ematching~can be seen as a hash join that only builds and
look-ups a single variable (i.e., $x$),
 and filters the outputs using the equality predicate on $\alpha$.
In other words, existing \ematching~algorithms
 will consider all $f(\alpha, g(\beta))$ terms regardless of whether $\alpha$ is congruent
 to $\beta$,
 while the generated conjunctive query gives the query optimizer the
 freedom to synthesize query plans that will consider only tuples where $\alpha\cong\beta$.

\subsection{Answering CQs with Generic Join}


% \begin{figure}
%   \includegraphics[width=0.9\linewidth]{trie.pdf}
%   \begin{minipage}{0.45\textwidth}
%     \subcaption{Trie for relation $R_f$ (\autoref{tab:rel-f}).}
%     \label{fig:trie-Rf}
%   \end{minipage}
%   \begin{minipage}{0.45\textwidth}
%     \subcaption{Trie for relation $R_g$ (\autoref{tab:rel-g}).}
%     \label{fig:trie-Rg}
%   \end{minipage}
%   \caption{
%     The tries built
%     for pattern $f(\alpha, g(\alpha))$
%     using variable ordering $[\alpha, x, \textit{root}]$
%     on \egraph from \autoref{fig:egraph-and-table}.
%   }\label{fig:trie}
% \end{figure}

Finally, we consider the problem of efficiently solving the compiled conjunctive
queries. We propose to use the generic join algorithm to solve the generated
conjunctive queries. Although traditional query plans, which are based on two
way joins such as hash joins and merge-sort joins, are extensively used in
industrial relational database engine, they may suffer on certain queries
compiled from patterns. For example, consider the pattern
$f(g(\alpha),h(\alpha))$. The compiled conjunctive query is:
\begin{equation}
  \label{eqn:cyclic-cq}
  Q(\textit{root},\alpha)\gets
  R_f(\textit{root}, x, y), R_g(x, \alpha), R_h(y, \alpha).
\end{equation}
Like the classic triangle query,
 this is a cyclic conjunctive query (\autoref{sec:hypergraph}).
We call \ematching~patterns that generate cyclic conjunctive
 queries \textit{cyclic patterns}.
For such cyclic queries,
 \citet{wcoj} show there exist databases on which {\em any}
 two-way join plan is suboptimal.
In contrast, generic join is guaranteed to run in time linear
 to the worst case output size.
Moreover, generic join can have comparable performance on acyclic
 queries with two-way join plans.
These properties make generic join our ideal
 solver for conjunctive queries generated from \ematching patterns.

% Suppose we fix the variable ordering $\pi$ to be $[x,y,\alpha,root]$.
% We first build a trie for $R_f$ on index
% $(\text{arg}_1,\text{arg}_2,\eclass)$ and a trie for $R_g$ on index
% $(\eclass, \text{arg}_1)$ as a preprocessing before actually  running the
% generic join
% algorithm. Then, the algorithm will intersect on
% column $\text{arg}_1$ of $R_f$ and column $\eclass$ of $R_g$, which nails down
% the possible values of $x$ and obtains residual relations $R_f[\text{arg}_1=x]$ and
% $R_g[\eclass=x]$. Next, it will intersect on column $\text{arg}_2$ of $R_f$ and
% column $\eclass$ of $R_g$ to nail down the values of $y$ and obtain residual
% relations $R_f[\text{arg}_1=x,\text{arg}_2=y]$ and $R_g[\eclass=y]$. To obtain a
% value for $\alpha$, the algorithm will then intersect on $R_g[\eclass=x]$ and
% $R_g[\eclass=y]$, and it will finally enumerate
% $R_f[\text{arg}_1=x,\text{arg}_2=y]$ to decide values for $\textit{root}$.

Using the generic join algorithm, suppose we fix the variable ordering to be
$[\alpha, x, \textit{root}]$ on the generated conjunctive query \ref{eqn:nonlinear-cq}.
The algorithm below shows generic join instantiated on this particular CQ:
% We first build
%  a trie for $R_f$ on index $(\mathit{child}_1, \mathit{child}_2, \mathit{root})$
%  a trie for $R_g$ on index $(\mathit{child}_1, \mathit{root})$
%  (tries shown in \autoref{fig:trie-Rf} and Figure~\ref{fig:trie-Rg}).
% Then, the algorithm will intersect on column $\mathit{child}_1$ of $R_f$ and on
% column $\mathit{child}_1$ of $R_g$, which nails down the possible values of
% $\alpha$ and obtains residual relation $R_f[\mathit{child}_1=\alpha]$ and
% $R_g[\mathit{child}_1=\alpha]$. Next, it will intersect on column $\mathit{child}_2$
% of $R_f[\mathit{child}_1=\alpha]$ and on $\eclass$ of
% $R_g[\mathit{child}_1=\alpha]$ to obtain possible values of $x$ and residual
% relation $R_f[\mathit{child}_1=\alpha, \mathit{child}_2=x]$. Finally, we enumerate
% terms of $R_f[\mathit{child}_1=\alpha, \mathit{child}_2=x]$ to find possible values
% of $\mathit{root}$.

\begin{algorithm}
  \LinesNumbered{} \SetAlgoLined{}
  \KwResult{compute $Q(\mathit{root},\alpha) \gets R_f(\mathit{root},\alpha,x),R_g(x,\alpha)$}
  \tcp{compute all possible values of $\alpha$}
  $A = R_f(\mathit{root},\alpha,x).\alpha \cap R_g(x,\alpha).\alpha$\;
  \For(){$i_\alpha \in A$} {
    \tcp{compute all possible values of $x$ given $\alpha=i_\alpha$}
    $X = R_f(\mathit{root},i_\alpha,x).x \cap R_g(x,i_\alpha).x$\;
    \For(){$i_x \in X$} {
      \tcp{compute all possible values of $\mathit{root}$ given $\alpha=i_\alpha$ and $x=i_x$}
      $\mathit{Roots} = R_f(\mathit{root},i_\alpha,i_x).\mathit{root} \cap R_g(i_x,i_\alpha).\mathit{root}$\;
      \For(){$i_{\mathit{root}} \in \mathit{Roots}$}{
        output$(i_{\mathit{root}}, i_\alpha)$ } } }
  \caption{
    Relational \ematching using GJ for $f(\alpha, g(\alpha))$,
    with ordering $[\alpha, x, \mathit{root}]$.
  }\label{alg:gj-pattern}
\end{algorithm}

\subsection{Complexity of Relational \Ematching}
% The translation of \ematching to query-answering on a database
%  is only productive if the size of the database is comparable to
%  that of the \egraph.
% Our relational encoding satisfies this requirement:
% \begin{lemma}\label{lem:bounded}
%   Let $I$ be the database representing the \egraph,
%    and let $N, C$ be the number of \enodes and \eclasses
%    in the \egraph, respectively.
%   We then have $\sum_{R \in I} |R| \leq N + C$,
%    where $|R| = \sum_{t \in R} |t|$
%    is the total size of the tuples in $R$,
%    and each value in the tuple has constant size.
% \end{lemma}

% \begin{proof}
%   Every \enode  corresponds to exactly one tuple,
%    and every tuple has exactly one extra attribute
%    than its \enode has children.
%   In addition, there may be \eclasses that are not children of
%    any \enodes.
%   Therefore $I$ has exactly $N$ tuples
%    and at most $N+C$ values.
% \end{proof}

Generic join guarantees worst-case optimality with respect to the output size,
 and relational \ematching preserves this optimality.
In particular, we have the following theorem:
\begin{theorem}
  Relational \ematching is worst-case optimal; that is, fix a pattern $p$, let
  $M(p,E)$ be the set of substitutions yielded by \ematching on an \egraph $E$
  with $N$ \enodes, relational \ematching runs in time $O(\max_E(|M(p,E)|))$.
\end{theorem}

\begin{proof}
  Notice that there is an one-to-one correspondence between output tuples of
  the generated conjunctive query and the \ematching pattern. Therefore, the
  worst-case bound is the same across an \ematching pattern and the conjunctive
  query it generated. Because generic join is worst-case optimal, relational
  \ematching also runs in worst-case optimal time with respect to the output
  size.
\end{proof}

The structure of \ematching patterns allows us to derive an additional bound
 dependent on the {\em actual} output size rather than the worst-case output size.
% Let $|\mathit{out}|$ denote the number of matches found,
%  we have the following:
\begin{theorem}
  Fix an \egraph $E$ with $N$ \enodes that compiles to a database $I$,
   and a fix pattern $p$ that compiles to conjunctive query
  $Q(\overline{X}) \gets R_1(\overline{X_1}),\ldots,R_m(\overline{X_m})$.
  Relational \ematching $p$ on $E$
   runs in time
   $O\left(\sqrt{|Q(I)| \times \Pi_{i}|R_i|}\right)
   \leq
    O\left(\sqrt{|Q(I)| \times N^m}\right)
   $.
% of a pattern with $|A|$ \enodes\
%    on an \egraph\ of size $N$ runs in time $O(\sqrt{N^m|Q(I)|})$.
\end{theorem}

\begin{proof}
  Let $\overline{X^\circ}$ be the set of isolated variables,
   those that occur in only one atom.
  Note that $\overline{X^\circ} \subseteq \overline{X}$,
   since $\overline{X}$ is precisely the pattern variables and the root,
   and auxiliary variables must occur in at least two atoms.
  Using these,
   define two new queries:
  \begin{align*}
  C(\overline{X^\circ}) & \gets R_{1}(\overline{X_{1}}), \ldots, R_m(\overline{X_m}) \\
  Q'(\overline{X}) & \gets R_{1}(\overline{X_{1}}), \ldots, R_m(\overline{X_m}), C(\overline{X^\circ})
  \end{align*}
 Since $\overline{X^\circ} \subseteq \overline{X}$,
  $C$ is the same query as $Q$ but with zero or more variables projected out.
 Therefore,
  every tuple in $C(I)$ corresponds to one in the output $Q(I)$,
  so $C(I) \subseteq Q(I)$ and $|C(I)| \leq |Q(I)|$.

 Now we can compute the AGM bound for $Q'$.
 Our new atom $C(\overline{X^\circ})$
  includes all those variables that only appear in one atom of $Q$.
 Therefore, every variable in $Q'$ occurs in at least two atoms,
  so assigning $1/2$ to each edge is a fractional edge cover.
 Thus:
 \begin{align*}
   \textsf{AGM}(Q') &= \sqrt{|C(I)| \times \Pi_i |R_i|} \\
                    &\leq \sqrt{|Q(I)| \times \Pi_i |R_i|} &\text{\quad since $|C(I)| \leq |Q(I)|$} \\
                    &\leq \sqrt{|Q(I)| \times N^m} &\text{\quad since $|R_i| < N$}
 \end{align*}

 Let $\GJ(Q', I)$ denote the running time
  of generic join with query $Q'$ on database $I$.
 We know that $\GJ(Q', I) \leq \textsf{AGM}(Q')$.
 Because $C(I) \subseteq Q(I)$, we also know $Q'(I) = Q(I)$,
  and we can use $GJ(Q', I)$ to bound $GJ(Q, I)$.
 Now we show that $\GJ(Q, I) \leq \GJ(Q', I)$.

 The query $Q'$ is just $Q$ with an additional
  atom $C(\overline{X^\circ})$ that
  covers the variables that only appeared in one atom from $Q$.
 Fix a variable ordering for generic join that puts those variables
  in $\overline{X^\circ}$ at the end.
 So loops of both \GJ instantiations are the same,
  except that, in $Q'$,
  each loop corresponding to a variable in $\overline{X^\circ}$ performs an intersection with $C$,
  but not in $Q$.
 But these intersections are in the innermost loops,
  at which point all intersections with atoms from $Q$ have already been done.
 So the intersections with $C$ do nothing,
  since $C$ is precisely $Q$ projected down to the variables in $\overline{X^\circ}$!
 Since those intersections are not helpful and $Q$ simply does not do them,
  $\GJ(Q, I) \leq \GJ(Q', I)$.

 Putting the inequalities together,
  we get:
 \[
   \GJ(Q, I)
   \leq
   \GJ(Q', I)
   \leq
   AGM(Q')
   \leq
   O\left(\sqrt{|Q(I)| \times \Pi_{i}|R_i|}\right)
   \leq
    O\left(\sqrt{|Q(I)| \times N^m}\right)
 \].
\end{proof}

\begin{example}[Complexity of relational \ematching]
  Consider the pattern $f(g(\alpha))$,
   which compiles to the query $Q(r, \alpha) \gets R_{f}(r, x), R_{g}(x, \alpha)$.
  Following the proof we define $C(r,\alpha) \gets R_{f}(r,x), R_{g}(x,\alpha)$
   and
   $Q'(r, \alpha) \gets R_{f}(r, x), R_{g}(x, \alpha), C(r, \alpha)$.
  The AGM bound for $Q'$ is $N^{1/2}N^{1/2}|C|^{1/2} = N\sqrt{|C|} = N\sqrt{|Q|}$.
  This also bounds the run time of generic join on $Q$.
\end{example}

The above bound is tight for linear patterns,
 in which case each variable occurs exactly twice in $Q'$.
In the case of nonlinear patterns,
 we may find tighter covers than assigning $1/2$ to each atom,
 thereby improving the bound.

\subsection{Supporting Multi-patterns}

Multi-patterns are an extension to \ematching used in both
 SMT solvers~\cite{efficient-ematching}
 and program optimizations~\cite{tensat}.
A multi-pattern is a list of patterns of the form $(p_1, \ldots, p_k)$
 that are to be simultaneously matched (i.e., the instantiation of each contained pattern
should use the same substitution $\sigma$).
For example, \ematching the multi-pattern
 $(f(\alpha, \beta), f(\alpha, \gamma))$ searches for pairs of two $f$-applications
 whose first arguments are equivalent.
Efficient support for multi-patterns
 on top of backtracking search requires complicated additions
 to state-of-the-art \ematching algorithms~\cite{efficient-ematching}.
Relational e-matching supports multi-patterns
 ``for free'':
 a multi-pattern is compiled to a single
 conjunctive query just like a single pattern.
For example, the conjunctive query for the
 multi-pattern $(f(\alpha, \beta), f(\alpha, \gamma))$ is
\begin{equation}
  Q(\textit{root}_1,\textit{root}_2,\alpha,\beta, \gamma)\gets
  R_f(\textit{root}_1,\alpha,\beta), R_g(\textit{root}_2, \alpha, \gamma).
\end{equation}
This is one example that shows the wide applicability of the relational
 model adopted in relational \ematching.


%%% Local Variables:
%%% mode: latex
%%% TeX-master: "main"
%%% End:
